\Chapter{115}

\Verse{1}
{
    \zR{话说宝玉为自己失言，被宝钗问住，想要掩饰过去，只见秋纹进来说：“外头 老爷叫二爷呢。”}
}
{
    \eR{[English source not present in the checked-in Bencraft/Joly files.]}
}
{
}

\Verse{2}
{
    \zR{宝玉巴不得一声儿，便走了。}
}
{
    \eR{[English source not present in the checked-in Bencraft/Joly files.]}
}
{
}

\Verse{3}
{
    \zR{到贾政那里，贾政道：“我叫你来不 为别的。}
}
{
    \eR{[English source not present in the checked-in Bencraft/Joly files.]}
}
{
}

\Verse{4}
{
    \zR{现在你穿着孝，不便到学里去，你在家里，必要将你念过的文章温习温习。}
}
{
    \eR{[English source not present in the checked-in Bencraft/Joly files.]}
}
{
}

\Verse{5}
{
    \zR{我这几天倒也闲着。}
}
{
    \eR{[English source not present in the checked-in Bencraft/Joly files.]}
}
{
}

\Verse{6}
{
    \zR{隔两三日要做几篇文章我瞧瞧，看你这些时进益了没有。”}
}
{
    \eR{[English source not present in the checked-in Bencraft/Joly files.]}
}
{
}

\Verse{7}
{
    \zR{宝 玉只得答应着。}
}
{
    \eR{[English source not present in the checked-in Bencraft/Joly files.]}
}
{
}

\Verse{8}
{
    \zR{贾政又道：“你环兄弟兰侄儿我也叫他们温习去了。}
}
{
    \eR{[English source not present in the checked-in Bencraft/Joly files.]}
}
{
}

\Verse{9}
{
    \zR{倘若你做的文 章不好，反倒不及他们，那可就不成事了。”}
}
{
    \eR{[English source not present in the checked-in Bencraft/Joly files.]}
}
{
}

\Verse{10}
{
    \zR{宝玉不敢言语，答应了个“是”，站着 不动。}
}
{
    \eR{[English source not present in the checked-in Bencraft/Joly files.]}
}
{
}

\Verse{11}
{
    \zR{贾政道：“去罢。”}
}
{
    \eR{[English source not present in the checked-in Bencraft/Joly files.]}
}
{
}

\Verse{12}
{
    \zR{宝玉退了出来，正遇见赖大诸人拿着些册子进来，宝玉一 溜烟回到自己房中。}
}
{
    \eR{[English source not present in the checked-in Bencraft/Joly files.]}
}
{
}

\Verse{13}
{
    \zR{宝钗问了，知道叫他作文章，倒也喜欢。}
}
{
    \eR{[English source not present in the checked-in Bencraft/Joly files.]}
}
{
}

\Verse{14}
{
    \zR{惟有宝玉不愿意，也 不敢怠慢。}
}
{
    \eR{[English source not present in the checked-in Bencraft/Joly files.]}
}
{
}

\Verse{15}
{
    \zR{正要坐下静静心，只见两个姑子进来，是地藏庵的。}
}
{
    \eR{[English source not present in the checked-in Bencraft/Joly files.]}
}
{
}

\Verse{16}
{
    \zR{见了宝钗，说道：“请二 奶奶安。”}
}
{
    \eR{[English source not present in the checked-in Bencraft/Joly files.]}
}
{
}

\Verse{17}
{
    \zR{宝钗待理不理的说：“你们好。”}
}
{
    \eR{[English source not present in the checked-in Bencraft/Joly files.]}
}
{
}

\Verse{18}
{
    \zR{因叫人来：“倒茶给师父们喝。”}
}
{
    \eR{[English source not present in the checked-in Bencraft/Joly files.]}
}
{
}

\Verse{19}
{
    \zR{宝玉原 要和那姑子说话，见宝钗似乎厌恶这些，也不好兜搭。}
}
{
    \eR{[English source not present in the checked-in Bencraft/Joly files.]}
}
{
}

\Verse{20}
{
    \zR{那姑子知道宝钗是个冷人， 也不久坐，辞了要去。}
}
{
    \eR{[English source not present in the checked-in Bencraft/Joly files.]}
}
{
}

\Verse{21}
{
    \zR{宝钗道：“再坐坐去罢。”}
}
{
    \eR{[English source not present in the checked-in Bencraft/Joly files.]}
}
{
}

\Verse{22}
{
    \zR{那姑子道：“我们因在铁槛寺做了 功德，好些时没来请太太奶奶们的安。}
}
{
    \eR{[English source not present in the checked-in Bencraft/Joly files.]}
}
{
}

\Verse{23}
{
    \zR{今日来了，见过了奶奶太太们，还要看看四 姑娘呢。”}
}
{
    \eR{[English source not present in the checked-in Bencraft/Joly files.]}
}
{
}

\Verse{24}
{
    \zR{宝钗点头，由他去了。}
}
{
    \eR{[English source not present in the checked-in Bencraft/Joly files.]}
}
{
}

\Verse{25}
{
    \zR{那姑子到了惜春那里，看见彩屏，便问：“姑娘在 那里呢？”}
}
{
    \eR{[English source not present in the checked-in Bencraft/Joly files.]}
}
{
}

\Verse{26}
{
    \zR{彩屏道：“不用提了。}
}
{
    \eR{[English source not present in the checked-in Bencraft/Joly files.]}
}
{
}

\Verse{27}
{
    \zR{姑娘这几天饭都没吃，只是歪着。”}
}
{
    \eR{[English source not present in the checked-in Bencraft/Joly files.]}
}
{
}

\Verse{28}
{
    \zR{那姑子道：“为 什么？”}
}
{
    \eR{[English source not present in the checked-in Bencraft/Joly files.]}
}
{
}

\Verse{29}
{
    \zR{彩屏道：“说也话长。}
}
{
    \eR{[English source not present in the checked-in Bencraft/Joly files.]}
}
{
}

\Verse{30}
{
    \zR{你见了姑娘，只怕他就和你说了。”}
}
{
    \eR{[English source not present in the checked-in Bencraft/Joly files.]}
}
{
}

\Verse{31}
{
    \zR{惜春早已听见， 急忙坐起，说：“你们两个人好啊，见我们家事差了，就不来了。”}
}
{
    \eR{[English source not present in the checked-in Bencraft/Joly files.]}
}
{
}

\Verse{32}
{
    \zR{那姑子道：“阿 弥陀佛!有也是施主，没也是施主，别说我们是本家庵里，受过老太太多少恩惠的。}
}
{
    \eR{[English source not present in the checked-in Bencraft/Joly files.]}
}
{
}

\Verse{33}
{
    \zR{如今老太太的事，太太奶奶们都见过了，只没有见姑娘，心里惦记，今儿是特特的 来瞧姑娘来了。”}
}
{
    \eR{[English source not present in the checked-in Bencraft/Joly files.]}
}
{
}

\Verse{34}
{
    \zR{惜春便问起水月庵的姑子来。}
}
{
    \eR{[English source not present in the checked-in Bencraft/Joly files.]}
}
{
}

\Verse{35}
{
    \zR{那姑子道：“他们庵里闹了些事，如今门上也不 肯常放进来了。”}
}
{
    \eR{[English source not present in the checked-in Bencraft/Joly files.]}
}
{
}

\Verse{36}
{
    \zR{便问惜春道：“前儿听见说，栊翠庵的妙师父怎么跟了人走了？”}
}
{
    \eR{[English source not present in the checked-in Bencraft/Joly files.]}
}
{
}

\Verse{37}
{
    \zR{惜春道：“那里的话？}
}
{
    \eR{[English source not present in the checked-in Bencraft/Joly files.]}
}
{
}

\Verse{38}
{
    \zR{说这个话的人提防着割舌头!人家遭了强盗抢去，怎么还说这 样的坏话。”}
}
{
    \eR{[English source not present in the checked-in Bencraft/Joly files.]}
}
{
}

\Verse{39}
{
    \zR{那姑子道：“妙师父的为人古怪，只怕是假惺惺罢？}
}
{
    \eR{[English source not present in the checked-in Bencraft/Joly files.]}
}
{
}

\Verse{40}
{
    \zR{在姑娘面前，我们 也不好说的。}
}
{
    \eR{[English source not present in the checked-in Bencraft/Joly files.]}
}
{
}

\Verse{41}
{
    \zR{那里像我们这些粗夯人，只知道讽经念佛，给人家忏悔，也为着自己 修个善果。”}
}
{
    \eR{[English source not present in the checked-in Bencraft/Joly files.]}
}
{
}

\Verse{42}
{
    \zR{惜春道：“怎么样就是善果呢？”}
}
{
    \eR{[English source not present in the checked-in Bencraft/Joly files.]}
}
{
}

\Verse{43}
{
    \zR{那姑子道：“除了咱们家这样善德人 家儿不怕，若是别人家那些诰命夫人小姐，也保不住一辈子的荣华。}
}
{
    \eR{[English source not present in the checked-in Bencraft/Joly files.]}
}
{
}

\Verse{44}
{
    \zR{到了苦难来了， 可就救不得了。}
}
{
    \eR{[English source not present in the checked-in Bencraft/Joly files.]}
}
{
}

\Verse{45}
{
    \zR{只有个观世音菩萨大慈大悲，遇见人家有苦难事，就慈心发动，设 法儿救济。}
}
{
    \eR{[English source not present in the checked-in Bencraft/Joly files.]}
}
{
}

\Verse{46}
{
    \zR{为什么如今都说‘大慈大悲救苦救难的观世音菩萨’呢。}
}
{
    \eR{[English source not present in the checked-in Bencraft/Joly files.]}
}
{
}

\Verse{47}
{
    \zR{我们修了行的 人，虽说比夫人小姐们苦多着呢，只是没有险难的了。}
}
{
    \eR{[English source not present in the checked-in Bencraft/Joly files.]}
}
{
}

\Verse{48}
{
    \zR{虽不能成佛作祖，修修来世 或者转个男身，自己也就好了。}
}
{
    \eR{[English source not present in the checked-in Bencraft/Joly files.]}
}
{
}

\Verse{49}
{
    \zR{不像如今脱生了个女人胎子，什么委屈烦难都说不 出来。}
}
{
    \eR{[English source not present in the checked-in Bencraft/Joly files.]}
}
{
}

\Verse{50}
{
    \zR{姑娘你还不知道呢，要是姑娘们到了出了门子，这一辈子跟着人，是更没法 儿的。}
}
{
    \eR{[English source not present in the checked-in Bencraft/Joly files.]}
}
{
}

\Verse{51}
{
    \zR{若说修行，也只要修得真。}
}
{
    \eR{[English source not present in the checked-in Bencraft/Joly files.]}
}
{
}

\Verse{52}
{
    \zR{那妙师父自为才情比我们强，他就嫌我们这些人 俗。}
}
{
    \eR{[English source not present in the checked-in Bencraft/Joly files.]}
}
{
}

\Verse{53}
{
    \zR{岂知俗的才能得善缘呢，他如今到底是遭了大劫了。”}
}
{
    \eR{[English source not present in the checked-in Bencraft/Joly files.]}
}
{
}

\Verse{54}
{
    \zR{惜春被那姑子一番话说的合在机上，也顾不得丫头们在这里，便将尤氏待他怎 样，前儿看家的事说了一遍，并将头发指给他瞧，道：“你打量我是什么没主意恋
        火坑的人么？}
}
{
    \eR{[English source not present in the checked-in Bencraft/Joly files.]}
}
{
}

\Verse{55}
{
    \zR{早有这样的心，只是想不出道儿来。”}
}
{
    \eR{[English source not present in the checked-in Bencraft/Joly files.]}
}
{
}

\Verse{56}
{
    \zR{那姑子听了，假作惊慌道：“姑 娘再别说这个话！}
}
{
    \eR{[English source not present in the checked-in Bencraft/Joly files.]}
}
{
}

\Verse{57}
{
    \zR{珍大奶奶听见，还要骂杀我们，撵出庵去呢。}
}
{
    \eR{[English source not present in the checked-in Bencraft/Joly files.]}
}
{
}

\Verse{58}
{
    \zR{姑娘这样人品，这 样人家，将来配个好姑爷，享一辈子的荣华富贵――”惜春不等说完，便红了脸， 说：“珍大奶奶撵得你，我就撵不得么？”}
}
{
    \eR{[English source not present in the checked-in Bencraft/Joly files.]}
}
{
}

\Verse{59}
{
    \zR{那姑子知是真心，便索性激他一激，说 道：“姑娘别怪我们说错了话。}
}
{
    \eR{[English source not present in the checked-in Bencraft/Joly files.]}
}
{
}

\Verse{60}
{
    \zR{太太奶奶们那里就依得姑娘的性子呢？}
}
{
    \eR{[English source not present in the checked-in Bencraft/Joly files.]}
}
{
}

\Verse{61}
{
    \zR{那时闹出没意 思来倒不好。}
}
{
    \eR{[English source not present in the checked-in Bencraft/Joly files.]}
}
{
}

\Verse{62}
{
    \zR{我们倒是为姑娘的话。”}
}
{
    \eR{[English source not present in the checked-in Bencraft/Joly files.]}
}
{
}

\Verse{63}
{
    \zR{惜春道：“这也瞧罢咧。”}
}
{
    \eR{[English source not present in the checked-in Bencraft/Joly files.]}
}
{
}

\Verse{64}
{
    \zR{彩屏等听这话头不 好，便使个眼色儿给姑子，叫他走。}
}
{
    \eR{[English source not present in the checked-in Bencraft/Joly files.]}
}
{
}

\Verse{65}
{
    \zR{那姑子会意，本来心里也害怕，不敢挑逗，便 告辞出去。}
}
{
    \eR{[English source not present in the checked-in Bencraft/Joly files.]}
}
{
}

\Verse{66}
{
    \zR{惜春也不留他，便冷笑道：“打量天下就是你们一个地藏庵么？”}
}
{
    \eR{[English source not present in the checked-in Bencraft/Joly files.]}
}
{
}

\Verse{67}
{
    \zR{那姑 子也不敢答言，去了。}
}
{
    \eR{[English source not present in the checked-in Bencraft/Joly files.]}
}
{
}

\Verse{68}
{
    \zR{彩屏见事不妥，恐耽不是，悄悄的去告诉了尤氏说：“四姑娘铰头发的念头还 没有息呢。}
}
{
    \eR{[English source not present in the checked-in Bencraft/Joly files.]}
}
{
}

\Verse{69}
{
    \zR{他这几天不是病，竟是怨命。}
}
{
    \eR{[English source not present in the checked-in Bencraft/Joly files.]}
}
{
}

\Verse{70}
{
    \zR{奶奶提防些，别闹出事来，那会子归罪我 们身上。”}
}
{
    \eR{[English source not present in the checked-in Bencraft/Joly files.]}
}
{
}

\Verse{71}
{
    \zR{尤氏道：“他那里是为要出家？}
}
{
    \eR{[English source not present in the checked-in Bencraft/Joly files.]}
}
{
}

\Verse{72}
{
    \zR{他为的是大爷不在家，安心和我过不去。}
}
{
    \eR{[English source not present in the checked-in Bencraft/Joly files.]}
}
{
}

\Verse{73}
{
    \zR{也只好由他罢了！”}
}
{
    \eR{[English source not present in the checked-in Bencraft/Joly files.]}
}
{
}

\Verse{74}
{
    \zR{彩屏等没法，也只好常常劝解。}
}
{
    \eR{[English source not present in the checked-in Bencraft/Joly files.]}
}
{
}

\Verse{75}
{
    \zR{岂知惜春一天一天的不吃饭， 只想铰头发。}
}
{
    \eR{[English source not present in the checked-in Bencraft/Joly files.]}
}
{
}

\Verse{76}
{
    \zR{彩屏等吃不住，只得到各处告诉。}
}
{
    \eR{[English source not present in the checked-in Bencraft/Joly files.]}
}
{
}

\Verse{77}
{
    \zR{邢王二夫人等也都劝了好几次，怎 奈惜春执迷不解。}
}
{
    \eR{[English source not present in the checked-in Bencraft/Joly files.]}
}
{
}

\Verse{78}
{
    \zR{邢王二夫人正要告诉贾政，只听外头传进来说：“甄家的太太带了他们家的宝 玉来了。”}
}
{
    \eR{[English source not present in the checked-in Bencraft/Joly files.]}
}
{
}

\Verse{79}
{
    \zR{众人急忙接出，便在王夫人处坐下。}
}
{
    \eR{[English source not present in the checked-in Bencraft/Joly files.]}
}
{
}

\Verse{80}
{
    \zR{众人行礼，叙些寒温，不必细述。}
}
{
    \eR{[English source not present in the checked-in Bencraft/Joly files.]}
}
{
}

\Verse{81}
{
    \zR{只言王夫人提起甄宝玉与自己的宝玉无二，要请甄宝玉进来一见。}
}
{
    \eR{[English source not present in the checked-in Bencraft/Joly files.]}
}
{
}

\Verse{82}
{
    \zR{传话出去，回来 说道：“甄少爷在外书房同老爷说话，说的投了机了，打发人来请我们二爷三爷， 还叫兰哥儿在外头吃饭，吃了饭进来。”}
}
{
    \eR{[English source not present in the checked-in Bencraft/Joly files.]}
}
{
}

\Verse{83}
{
    \zR{说毕，里头也便摆饭。}
}
{
    \eR{[English source not present in the checked-in Bencraft/Joly files.]}
}
{
}

\Verse{84}
{
    \zR{原来此时贾政见甄宝玉相貌果与宝玉一样，试探他的文才，竟应对如流，甚是 心敬，故叫宝玉等三人出来警励他们，再者到底叫宝玉来比一比。}
}
{
    \eR{[English source not present in the checked-in Bencraft/Joly files.]}
}
{
}

\Verse{85}
{
    \zR{宝玉听命，穿了 素服，带了兄弟侄儿出来，见了甄宝玉，竟是旧相识一般。}
}
{
    \eR{[English source not present in the checked-in Bencraft/Joly files.]}
}
{
}

\Verse{86}
{
    \zR{那甄宝玉也像那里见过 的。}
}
{
    \eR{[English source not present in the checked-in Bencraft/Joly files.]}
}
{
}

\Verse{87}
{
    \zR{两人行了礼，然后贾环贾兰相见。}
}
{
    \eR{[English source not present in the checked-in Bencraft/Joly files.]}
}
{
}

\Verse{88}
{
    \zR{本来贾政席地而坐，要让甄宝玉在椅子上坐， 甄宝玉因是晚辈，不敢上坐，就在地下铺了褥子坐下。}
}
{
    \eR{[English source not present in the checked-in Bencraft/Joly files.]}
}
{
}

\Verse{89}
{
    \zR{如今宝玉等出来，又不能同 贾政一处坐着，为甄宝玉是晚一辈，又不好竟叫宝玉等站着。}
}
{
    \eR{[English source not present in the checked-in Bencraft/Joly files.]}
}
{
}

\Verse{90}
{
    \zR{贾政知是不便，站起 来又说了几句话，叫人摆饭，说：“我失陪，叫小儿辈陪着，大家说话儿，好叫他 们领领大教。”}
}
{
    \eR{[English source not present in the checked-in Bencraft/Joly files.]}
}
{
}

\Verse{91}
{
    \zR{甄宝玉逊谢道：“老伯大人请便，小侄正欲领世兄们的教呢。”}
}
{
    \eR{[English source not present in the checked-in Bencraft/Joly files.]}
}
{
}

\Verse{92}
{
    \zR{贾政 回复了几句，便自往内书房去。}
}
{
    \eR{[English source not present in the checked-in Bencraft/Joly files.]}
}
{
}

\Verse{93}
{
    \zR{那甄宝玉却要送出来，贾政拦住。}
}
{
    \eR{[English source not present in the checked-in Bencraft/Joly files.]}
}
{
}

\Verse{94}
{
    \zR{宝玉等先抢了一 步，出了书房门槛站立着，看贾政进去，然后进来让甄宝玉坐下。}
}
{
    \eR{[English source not present in the checked-in Bencraft/Joly files.]}
}
{
}

\Verse{95}
{
    \zR{彼此套叙了一回， 诸如久慕渴想的话，也不必细述。}
}
{
    \eR{[English source not present in the checked-in Bencraft/Joly files.]}
}
{
}

\Verse{96}
{
    \zR{且说贾宝玉见了甄宝玉，想到梦中之景，并且素知甄宝玉为人，必是和他同心， 以为得了知己。}
}
{
    \eR{[English source not present in the checked-in Bencraft/Joly files.]}
}
{
}

\Verse{97}
{
    \zR{因初次见面，不便造次，且又贾环贾兰在坐，只有极力夸赞说：“久 仰芳名，无由亲炙，今日见面，真是谪仙一流的人物。”}
}
{
    \eR{[English source not present in the checked-in Bencraft/Joly files.]}
}
{
}

\Verse{98}
{
    \zR{那甄宝玉素来也知贾宝玉 的为人，今日一见，果然不差，“只是可与我共学，不可与我适道。}
}
{
    \eR{[English source not present in the checked-in Bencraft/Joly files.]}
}
{
}

\Verse{99}
{
    \zR{他既和我同名 同貌，也是三生石上的旧精魂了。}
}
{
    \eR{[English source not present in the checked-in Bencraft/Joly files.]}
}
{
}

\Verse{100}
{
    \zR{我如今略知些道理，何不和他讲讲？}
}
{
    \eR{[English source not present in the checked-in Bencraft/Joly files.]}
}
{
}

\Verse{101}
{
    \zR{但只是初见， 尚不知他的心与我同不同，只好缓缓的来。”}
}
{
    \eR{[English source not present in the checked-in Bencraft/Joly files.]}
}
{
}

\Verse{102}
{
    \zR{便道：“世兄的才名，弟所素知的。}
}
{
    \eR{[English source not present in the checked-in Bencraft/Joly files.]}
}
{
}

\Verse{103}
{
    \zR{在 世兄是数万人里头选出来最清最雅的。}
}
{
    \eR{[English source not present in the checked-in Bencraft/Joly files.]}
}
{
}

\Verse{104}
{
    \zR{至于弟乃庸庸碌碌一等愚人，忝附同名，殊 觉玷辱了这两个字。”}
}
{
    \eR{[English source not present in the checked-in Bencraft/Joly files.]}
}
{
}

\Verse{105}
{
    \zR{贾宝玉听了，心想：“这个人果然同我的心一样的，但是你我 都是男人，不比那女孩儿们清洁，怎么他拿我当作女孩儿看待起来？”}
}
{
    \eR{[English source not present in the checked-in Bencraft/Joly files.]}
}
{
}

\Verse{106}
{
    \zR{便道：“世 兄谬赞，实不敢当。}
}
{
    \eR{[English source not present in the checked-in Bencraft/Joly files.]}
}
{
}

\Verse{107}
{
    \zR{弟至浊至愚，只不过一块顽石耳，何敢比世兄品望清高，实称 此两字呢？”}
}
{
    \eR{[English source not present in the checked-in Bencraft/Joly files.]}
}
{
}

\Verse{108}
{
    \zR{甄宝玉道：“弟少时不知分量，自谓尚可琢磨；岂知家遭消索，数年 来更比瓦砾犹贱。}
}
{
    \eR{[English source not present in the checked-in Bencraft/Joly files.]}
}
{
}

\Verse{109}
{
    \zR{虽不敢说历尽甘苦，然世道人情，略略的领悟了些须。}
}
{
    \eR{[English source not present in the checked-in Bencraft/Joly files.]}
}
{
}

\Verse{110}
{
    \zR{世兄是锦 衣玉食，无不遂心的，必是文章经济高出人上，所以老伯钟爱，将为席上之珍。}
}
{
    \eR{[English source not present in the checked-in Bencraft/Joly files.]}
}
{
}

\Verse{111}
{
    \zR{弟 所以才说尊名方称。”}
}
{
    \eR{[English source not present in the checked-in Bencraft/Joly files.]}
}
{
}

\Verse{112}
{
    \zR{贾宝玉听这话头又近了禄蠹的旧套，想话回答。}
}
{
    \eR{[English source not present in the checked-in Bencraft/Joly files.]}
}
{
}

\Verse{113}
{
    \zR{贾环见未与 他说话，心中早不自在。}
}
{
    \eR{[English source not present in the checked-in Bencraft/Joly files.]}
}
{
}

\Verse{114}
{
    \zR{倒是贾兰听了这话，甚觉合意，便说道：“世叔所言，固 是太谦，若论到文章经济，实在从历练中出来的，方为真才实学。}
}
{
    \eR{[English source not present in the checked-in Bencraft/Joly files.]}
}
{
}

\Verse{115}
{
    \zR{在小侄年幼，虽 不知文章为何物，然将读过的细味起来，那膏粱文绣，比着令闻广誉，真是不啻百 倍的了！”}
}
{
    \eR{[English source not present in the checked-in Bencraft/Joly files.]}
}
{
}

\Verse{116}
{
    \zR{甄宝玉未及答言。}
}
{
    \eR{[English source not present in the checked-in Bencraft/Joly files.]}
}
{
}

\Verse{117}
{
    \zR{贾宝玉听了兰儿的话，心里越发不合，想道：“这孩子从几时也学了这一派酸 论！”}
}
{
    \eR{[English source not present in the checked-in Bencraft/Joly files.]}
}
{
}

\Verse{118}
{
    \zR{便说道：“弟闻得世兄也诋尽流俗，性情中另有一番见解。}
}
{
    \eR{[English source not present in the checked-in Bencraft/Joly files.]}
}
{
}

\Verse{119}
{
    \zR{今日弟幸会芝范， 想欲领教一番超凡入圣的道理，从此可以洗净俗肠，重开眼界。}
}
{
    \eR{[English source not present in the checked-in Bencraft/Joly files.]}
}
{
}

\Verse{120}
{
    \zR{不意视弟为蠢物， 所以将世路的话来酬应。”}
}
{
    \eR{[English source not present in the checked-in Bencraft/Joly files.]}
}
{
}

\Verse{121}
{
    \zR{甄宝玉听说，心里晓得：“他知我少年的性情，所以疑我 为假。}
}
{
    \eR{[English source not present in the checked-in Bencraft/Joly files.]}
}
{
}

\Verse{122}
{
    \zR{我索性把话说明，或者与我作个知心朋友，也是好的。”}
}
{
    \eR{[English source not present in the checked-in Bencraft/Joly files.]}
}
{
}

\Verse{123}
{
    \zR{便说：“世兄高论， 固是真切。}
}
{
    \eR{[English source not present in the checked-in Bencraft/Joly files.]}
}
{
}

\Verse{124}
{
    \zR{但弟少时也曾深恶那些旧套陈言，只是一年长似一年，家君致仕在家， 懒于酬应，委弟接待。}
}
{
    \eR{[English source not present in the checked-in Bencraft/Joly files.]}
}
{
}

\Verse{125}
{
    \zR{后来见过那些大人先生，尽都是显亲扬名的人；便是著书立 说，无非言忠言孝，自有一番立德立言的事业，方不枉生在圣明之时，也不致负了 父亲师长养育教诲之恩。}
}
{
    \eR{[English source not present in the checked-in Bencraft/Joly files.]}
}
{
}

\Verse{126}
{
    \zR{所以把少时那些迂想痴情，渐渐的淘汰了些。}
}
{
    \eR{[English source not present in the checked-in Bencraft/Joly files.]}
}
{
}

\Verse{127}
{
    \zR{如今尚欲访 师觅友，教导愚蒙。}
}
{
    \eR{[English source not present in the checked-in Bencraft/Joly files.]}
}
{
}

\Verse{128}
{
    \zR{幸会世兄，定当有以教我。}
}
{
    \eR{[English source not present in the checked-in Bencraft/Joly files.]}
}
{
}

\Verse{129}
{
    \zR{适才所言，并非虚意。”}
}
{
    \eR{[English source not present in the checked-in Bencraft/Joly files.]}
}
{
}

\Verse{130}
{
    \zR{贾宝玉愈 听愈不耐烦，又不好冷淡，只得将言语支吾。}
}
{
    \eR{[English source not present in the checked-in Bencraft/Joly files.]}
}
{
}

\Verse{131}
{
    \zR{幸喜里头传出话来，说：“若是外头 爷们吃了饭，请甄少爷里头去坐呢。”}
}
{
    \eR{[English source not present in the checked-in Bencraft/Joly files.]}
}
{
}

\Verse{132}
{
    \zR{宝玉听了，趁势便邀甄宝玉进去。}
}
{
    \eR{[English source not present in the checked-in Bencraft/Joly files.]}
}
{
}

\Verse{133}
{
    \zR{那甄宝玉 依命前行，贾宝玉等陪着来见王夫人。}
}
{
    \eR{[English source not present in the checked-in Bencraft/Joly files.]}
}
{
}

\Verse{134}
{
    \zR{贾宝玉见是甄太太上坐，便先请过了安。}
}
{
    \eR{[English source not present in the checked-in Bencraft/Joly files.]}
}
{
}

\Verse{135}
{
    \zR{贾 环贾兰也见了。}
}
{
    \eR{[English source not present in the checked-in Bencraft/Joly files.]}
}
{
}

\Verse{136}
{
    \zR{甄宝玉也请了王夫人的安。}
}
{
    \eR{[English source not present in the checked-in Bencraft/Joly files.]}
}
{
}

\Verse{137}
{
    \zR{两母两子，互相厮认。}
}
{
    \eR{[English source not present in the checked-in Bencraft/Joly files.]}
}
{
}

\Verse{138}
{
    \zR{虽是贾宝玉是娶 过亲的，那甄夫人年纪已老，又是老亲，因见贾宝玉的相貌身材与他儿子一般，不 禁亲热起来。}
}
{
    \eR{[English source not present in the checked-in Bencraft/Joly files.]}
}
{
}

\Verse{139}
{
    \zR{王夫人更不用说，拉着甄宝玉问长问短，觉得比自己家的宝玉老成些。}
}
{
    \eR{[English source not present in the checked-in Bencraft/Joly files.]}
}
{
}

\Verse{140}
{
    \zR{回看贾兰，也是清秀超群的，虽不能像两个宝玉的形象，也还随得上，只有贾环粗 夯，未免有偏爱之色。}
}
{
    \eR{[English source not present in the checked-in Bencraft/Joly files.]}
}
{
}

\Verse{141}
{
    \zR{众人一见两个宝玉在这里，都来瞧看，说道：“真真奇事!名字同了也罢，怎么 相貌身材都是一样的。}
}
{
    \eR{[English source not present in the checked-in Bencraft/Joly files.]}
}
{
}

\Verse{142}
{
    \zR{亏得是我们宝玉穿孝，若是一样的衣服穿着，一时也认不出 来。”}
}
{
    \eR{[English source not present in the checked-in Bencraft/Joly files.]}
}
{
}

\Verse{143}
{
    \zR{内中紫鹃一时痴意发作，因想起黛玉来，心里说道：“可惜林姑娘死了，若不 死时，就将那甄宝玉配了他，只怕也是愿意的。”}
}
{
    \eR{[English source not present in the checked-in Bencraft/Joly files.]}
}
{
}

\Verse{144}
{
    \zR{正想着，只听得甄夫人道：“前日 听得我们老爷回来说：我们宝玉年纪也大了，求这里老爷留心一门亲事。”}
}
{
    \eR{[English source not present in the checked-in Bencraft/Joly files.]}
}
{
}

\Verse{145}
{
    \zR{王夫人 正爱甄宝玉，顺口便说道：“我也想要与令郎作伐。}
}
{
    \eR{[English source not present in the checked-in Bencraft/Joly files.]}
}
{
}

\Verse{146}
{
    \zR{我家有四个姑娘：那三个都不 用说，死的死，嫁的嫁了。}
}
{
    \eR{[English source not present in the checked-in Bencraft/Joly files.]}
}
{
}

\Verse{147}
{
    \zR{还有我们珍大侄儿的妹子，只是年纪过小几岁，恐怕难 配。}
}
{
    \eR{[English source not present in the checked-in Bencraft/Joly files.]}
}
{
}

\Verse{148}
{
    \zR{倒是我们大媳妇的两个堂妹子，生得人材齐正。}
}
{
    \eR{[English source not present in the checked-in Bencraft/Joly files.]}
}
{
}

\Verse{149}
{
    \zR{二姑娘呢，已经许了人家；三 姑娘正好与令郎为配。}
}
{
    \eR{[English source not present in the checked-in Bencraft/Joly files.]}
}
{
}

\Verse{150}
{
    \zR{过一天，我给令郎作媒。}
}
{
    \eR{[English source not present in the checked-in Bencraft/Joly files.]}
}
{
}

\Verse{151}
{
    \zR{但是他家的家计如今差些。”}
}
{
    \eR{[English source not present in the checked-in Bencraft/Joly files.]}
}
{
}

\Verse{152}
{
    \zR{甄夫 人道：“太太这话又客套了。}
}
{
    \eR{[English source not present in the checked-in Bencraft/Joly files.]}
}
{
}

\Verse{153}
{
    \zR{如今我们家还有什么？}
}
{
    \eR{[English source not present in the checked-in Bencraft/Joly files.]}
}
{
}

\Verse{154}
{
    \zR{只怕人家嫌我们穷罢咧。”}
}
{
    \eR{[English source not present in the checked-in Bencraft/Joly files.]}
}
{
}

\Verse{155}
{
    \zR{王夫 人道：“现今府上复又出了差，将来不但复旧，必是比先前更要鼎盛起来。”}
}
{
    \eR{[English source not present in the checked-in Bencraft/Joly files.]}
}
{
}

\Verse{156}
{
    \zR{甄夫人 笑着道：“但愿依着太太的话更好。}
}
{
    \eR{[English source not present in the checked-in Bencraft/Joly files.]}
}
{
}

\Verse{157}
{
    \zR{这么着，就求太太作个保山。”}
}
{
    \eR{[English source not present in the checked-in Bencraft/Joly files.]}
}
{
}

\Verse{158}
{
    \zR{甄宝玉听见他们 说起亲事，便告辞出来，贾宝玉等只得陪着来到书房。}
}
{
    \eR{[English source not present in the checked-in Bencraft/Joly files.]}
}
{
}

\Verse{159}
{
    \zR{见贾政已在那里，复又立谈 几句。}
}
{
    \eR{[English source not present in the checked-in Bencraft/Joly files.]}
}
{
}

\Verse{160}
{
    \zR{听见甄家的人来回甄宝玉道：“太太要走了，请爷回去罢。”}
}
{
    \eR{[English source not present in the checked-in Bencraft/Joly files.]}
}
{
}

\Verse{161}
{
    \zR{于是甄宝玉告辞 出来。}
}
{
    \eR{[English source not present in the checked-in Bencraft/Joly files.]}
}
{
}

\Verse{162}
{
    \zR{贾政命宝玉、环、兰相送，不提。}
}
{
    \eR{[English source not present in the checked-in Bencraft/Joly files.]}
}
{
}

\Verse{163}
{
    \zR{且说宝玉自那日见了甄宝玉之父，知道甄宝玉来京，朝夕盼望。}
}
{
    \eR{[English source not present in the checked-in Bencraft/Joly files.]}
}
{
}

\Verse{164}
{
    \zR{今儿见面，原 想得一知己，岂知谈了半天，竟有些冰炭不投。}
}
{
    \eR{[English source not present in the checked-in Bencraft/Joly files.]}
}
{
}

\Verse{165}
{
    \zR{闷闷的回到自己房中，也不言，也 不笑，只管发怔。}
}
{
    \eR{[English source not present in the checked-in Bencraft/Joly files.]}
}
{
}

\Verse{166}
{
    \zR{宝钗便问：“那甄宝玉果然像你么？”}
}
{
    \eR{[English source not present in the checked-in Bencraft/Joly files.]}
}
{
}

\Verse{167}
{
    \zR{宝玉道：“相貌倒还是一样 的，只是言谈间看起来，并不知道什么，不过也是个禄蠹。”}
}
{
    \eR{[English source not present in the checked-in Bencraft/Joly files.]}
}
{
}

\Verse{168}
{
    \zR{宝钗道：“你又编派人 家了。}
}
{
    \eR{[English source not present in the checked-in Bencraft/Joly files.]}
}
{
}

\Verse{169}
{
    \zR{怎么就见得也是个禄蠹呢？”}
}
{
    \eR{[English source not present in the checked-in Bencraft/Joly files.]}
}
{
}

\Verse{170}
{
    \zR{宝玉道：“他说了半天，并没个明心见性之谈， 不过说些什么‘文章经济’，又说什么‘为忠为孝’。}
}
{
    \eR{[English source not present in the checked-in Bencraft/Joly files.]}
}
{
}

\Verse{171}
{
    \zR{这样人可不是个禄蠹么？}
}
{
    \eR{[English source not present in the checked-in Bencraft/Joly files.]}
}
{
}

\Verse{172}
{
    \zR{只可 惜他也生了这样一个相貌。}
}
{
    \eR{[English source not present in the checked-in Bencraft/Joly files.]}
}
{
}

\Verse{173}
{
    \zR{我想来，有了他，我竟要连我这个相貌都不要了。”}
}
{
    \eR{[English source not present in the checked-in Bencraft/Joly files.]}
}
{
}

\Verse{174}
{
    \zR{宝 钗见他又说呆话，便说道：“你真真说出句话来叫人发笑，这相貌怎么能不要呢!况 且人家这话是正理，做了一个男人，原该要立身扬名的，谁像你一味的柔情私意？}
}
{
    \eR{[English source not present in the checked-in Bencraft/Joly files.]}
}
{
}

\Verse{175}
{
    \zR{不说自己没有刚烈，倒说人家是禄蠹。”}
}
{
    \eR{[English source not present in the checked-in Bencraft/Joly files.]}
}
{
}

\Verse{176}
{
    \zR{宝玉本听了甄宝玉的话，甚不耐烦，又被 宝钗抢白了一场，心中更加不乐，闷闷昏昏，不觉将旧病又勾起来了，并不言语， 只是傻笑。}
}
{
    \eR{[English source not present in the checked-in Bencraft/Joly files.]}
}
{
}

\Verse{177}
{
    \zR{宝钗不知，只道自己的话错了，他所以冷笑，也不理他。}
}
{
    \eR{[English source not present in the checked-in Bencraft/Joly files.]}
}
{
}

\Verse{178}
{
    \zR{岂知那日便有 些发呆，袭人等怄他，也不言语。}
}
{
    \eR{[English source not present in the checked-in Bencraft/Joly files.]}
}
{
}

\Verse{179}
{
    \zR{过了一夜，次日起来，只是呆呆的，竟有前番病 的样子。}
}
{
    \eR{[English source not present in the checked-in Bencraft/Joly files.]}
}
{
}

\Verse{180}
{
    \zR{一日，王夫人因为惜春定要铰发出家，尤氏不能拦阻，看着惜春的样子是若不 依他必要自尽的，虽然昼夜着人看守终非常事，便告诉了贾政。}
}
{
    \eR{[English source not present in the checked-in Bencraft/Joly files.]}
}
{
}

\Verse{181}
{
    \zR{贾政叹气跺脚，只 说：“东府里不知干了什么，闹到如此地位！”}
}
{
    \eR{[English source not present in the checked-in Bencraft/Joly files.]}
}
{
}

\Verse{182}
{
    \zR{叫了贾蓉来说了一顿，叫他去和他母 亲说：“认真劝解劝解。}
}
{
    \eR{[English source not present in the checked-in Bencraft/Joly files.]}
}
{
}

\Verse{183}
{
    \zR{若是必要这样，就不是我们家的姑娘了。”}
}
{
    \eR{[English source not present in the checked-in Bencraft/Joly files.]}
}
{
}

\Verse{184}
{
    \zR{岂知尤氏不劝还 好，一劝了，更要寻死，说：“做了女孩儿，终不能在家一辈子的。}
}
{
    \eR{[English source not present in the checked-in Bencraft/Joly files.]}
}
{
}

\Verse{185}
{
    \zR{若像二姐姐一 样，老爷太太们倒要操心，况且死了。}
}
{
    \eR{[English source not present in the checked-in Bencraft/Joly files.]}
}
{
}

\Verse{186}
{
    \zR{如今譬如我死了似的，放我出了家，干干净 净的一辈子，就是疼我了。}
}
{
    \eR{[English source not present in the checked-in Bencraft/Joly files.]}
}
{
}

\Verse{187}
{
    \zR{况且我又不出门，就是栊翠庵原是咱们家的基址，我就 在那里修行。}
}
{
    \eR{[English source not present in the checked-in Bencraft/Joly files.]}
}
{
}

\Verse{188}
{
    \zR{我有什么，你们也照应得着。}
}
{
    \eR{[English source not present in the checked-in Bencraft/Joly files.]}
}
{
}

\Verse{189}
{
    \zR{现在妙玉的当家的在那里。}
}
{
    \eR{[English source not present in the checked-in Bencraft/Joly files.]}
}
{
}

\Verse{190}
{
    \zR{你们依我呢， 我就算得了命了；若不依我呢，我也没法，只有死就完了!我如若遂了自己的心愿，
        那时哥哥回来，我和他说并不是你们逼着我的；若说我死了，未免哥哥回来，倒说 你们不容我。”}
}
{
    \eR{[English source not present in the checked-in Bencraft/Joly files.]}
}
{
}

\Verse{191}
{
    \zR{尤氏本与惜春不合，听他的话，也似乎有理，只得去回王夫人。}
}
{
    \eR{[English source not present in the checked-in Bencraft/Joly files.]}
}
{
}

\Verse{192}
{
    \zR{王夫人已到宝钗那里，见宝玉神魂失所，心下着忙，便说袭人道：“你们忒不 留神!二爷犯了病，也不来回我。”}
}
{
    \eR{[English source not present in the checked-in Bencraft/Joly files.]}
}
{
}

\Verse{193}
{
    \zR{袭人道：“二爷的病原来是常有的，一时好，一 时不好。}
}
{
    \eR{[English source not present in the checked-in Bencraft/Joly files.]}
}
{
}

\Verse{194}
{
    \zR{天天到太太那里，仍旧请安去，原是好好儿的，今日才发糊涂些。}
}
{
    \eR{[English source not present in the checked-in Bencraft/Joly files.]}
}
{
}

\Verse{195}
{
    \zR{二奶奶 正要来回太太，恐怕太太说我们大惊小怪。”}
}
{
    \eR{[English source not present in the checked-in Bencraft/Joly files.]}
}
{
}

\Verse{196}
{
    \zR{宝玉听见王夫人说他们，心里一时明 白，怕他们受委屈，便说道：“太太放心，我没什么病，只是心里觉着有些闷闷的。”}
}
{
    \eR{[English source not present in the checked-in Bencraft/Joly files.]}
}
{
}

\Verse{197}
{
    \zR{王夫人道：“你是有这病根子，早说了，好请大夫瞧瞧，吃两剂药好了不好？}
}
{
    \eR{[English source not present in the checked-in Bencraft/Joly files.]}
}
{
}

\Verse{198}
{
    \zR{若再闹 到头里丢了玉的样子，那可就费事了。”}
}
{
    \eR{[English source not present in the checked-in Bencraft/Joly files.]}
}
{
}

\Verse{199}
{
    \zR{宝玉道：“太太不放心，便叫个人瞧瞧，我 就吃药。”}
}
{
    \eR{[English source not present in the checked-in Bencraft/Joly files.]}
}
{
}

\Verse{200}
{
    \zR{王夫人便叫丫头传话出来请大夫。}
}
{
    \eR{[English source not present in the checked-in Bencraft/Joly files.]}
}
{
}

\Verse{201}
{
    \zR{这一个心思都在宝玉身上，便将惜春 的事忘了。}
}
{
    \eR{[English source not present in the checked-in Bencraft/Joly files.]}
}
{
}

\Verse{202}
{
    \zR{迟了一回，大夫看了服药，王夫人回去。}
}
{
    \eR{[English source not present in the checked-in Bencraft/Joly files.]}
}
{
}

\Verse{203}
{
    \zR{过了几天，宝玉更糊涂了，甚至于饮食不进，大家着急起来。}
}
{
    \eR{[English source not present in the checked-in Bencraft/Joly files.]}
}
{
}

\Verse{204}
{
    \zR{恰又忙着脱孝， 家中无人，又叫了贾芸来照应大夫。}
}
{
    \eR{[English source not present in the checked-in Bencraft/Joly files.]}
}
{
}

\Verse{205}
{
    \zR{贾琏家下无人，请了王仁来在外帮着料理。}
}
{
    \eR{[English source not present in the checked-in Bencraft/Joly files.]}
}
{
}

\Verse{206}
{
    \zR{那 巧姐儿是日夜哭母，也是病了。}
}
{
    \eR{[English source not present in the checked-in Bencraft/Joly files.]}
}
{
}

\Verse{207}
{
    \zR{所以荣府中又闹得马仰人翻。}
}
{
    \eR{[English source not present in the checked-in Bencraft/Joly files.]}
}
{
}

\Verse{208}
{
    \zR{一日，又当脱孝来家，王夫人亲身又看宝玉，见宝玉人事不醒，急得众人手足 无措，一面哭着，一面告诉贾政说：“大夫说了，不肯下药，只好预备后事！”}
}
{
    \eR{[English source not present in the checked-in Bencraft/Joly files.]}
}
{
}

\Verse{209}
{
    \zR{贾政 叹气连连，只得亲自看视，见其光景果然不好，便又叫贾琏办去。}
}
{
    \eR{[English source not present in the checked-in Bencraft/Joly files.]}
}
{
}

\Verse{210}
{
    \zR{贾琏不敢违拗， 只得叫人料理；手头又短，正在为难。}
}
{
    \eR{[English source not present in the checked-in Bencraft/Joly files.]}
}
{
}

\Verse{211}
{
    \zR{只见一个人跳进来说：“二爷不好了，又有 饥荒来了！”}
}
{
    \eR{[English source not present in the checked-in Bencraft/Joly files.]}
}
{
}

\Verse{212}
{
    \zR{贾琏不知何事，这一吓非同小可，瞪着眼说道：“什么事？”}
}
{
    \eR{[English source not present in the checked-in Bencraft/Joly files.]}
}
{
}

\Verse{213}
{
    \zR{那小厮道： “门上来了一个和尚，手里拿着二爷的这块丢的玉，说要一万赏银。”}
}
{
    \eR{[English source not present in the checked-in Bencraft/Joly files.]}
}
{
}

\Verse{214}
{
    \zR{贾琏照脸啐 道：“我打量什么事，这样慌张!前番那假的你不知道么？}
}
{
    \eR{[English source not present in the checked-in Bencraft/Joly files.]}
}
{
}

\Verse{215}
{
    \zR{就是真的，现在人要死了， 要这玉做什么？”}
}
{
    \eR{[English source not present in the checked-in Bencraft/Joly files.]}
}
{
}

\Verse{216}
{
    \zR{小厮道：“奴才也说了。}
}
{
    \eR{[English source not present in the checked-in Bencraft/Joly files.]}
}
{
}

\Verse{217}
{
    \zR{那和尚说，给他银子就好了。”}
}
{
    \eR{[English source not present in the checked-in Bencraft/Joly files.]}
}
{
}

\Verse{218}
{
    \zR{正说着， 外头嚷进来说：“这和尚撒野，各自跑进来了，众人拦他拦不住！”}
}
{
    \eR{[English source not present in the checked-in Bencraft/Joly files.]}
}
{
}

\Verse{219}
{
    \zR{贾琏道：“那里 有这样怪事？}
}
{
    \eR{[English source not present in the checked-in Bencraft/Joly files.]}
}
{
}

\Verse{220}
{
    \zR{你们还不快打出去呢。”}
}
{
    \eR{[English source not present in the checked-in Bencraft/Joly files.]}
}
{
}

\Verse{221}
{
    \zR{又闹着，贾政听见了，也没了主意了。}
}
{
    \eR{[English source not present in the checked-in Bencraft/Joly files.]}
}
{
}

\Verse{222}
{
    \zR{里头又 哭出来，说：“宝二爷不好了！”}
}
{
    \eR{[English source not present in the checked-in Bencraft/Joly files.]}
}
{
}

\Verse{223}
{
    \zR{贾政益发着急。}
}
{
    \eR{[English source not present in the checked-in Bencraft/Joly files.]}
}
{
}

\Verse{224}
{
    \zR{只见那和尚说道：“要命拿银子来。”}
}
{
    \eR{[English source not present in the checked-in Bencraft/Joly files.]}
}
{
}

\Verse{225}
{
    \zR{贾政忽然想起：“头里宝玉的病是和尚治好的；这会子和尚来，或者有救星。}
}
{
    \eR{[English source not present in the checked-in Bencraft/Joly files.]}
}
{
}

\Verse{226}
{
    \zR{但是 这玉倘或是真，他要起银子来，怎么样呢？”}
}
{
    \eR{[English source not present in the checked-in Bencraft/Joly files.]}
}
{
}

\Verse{227}
{
    \zR{想一想：“如今且不管他，果真人好 了再说。”}
}
{
    \eR{[English source not present in the checked-in Bencraft/Joly files.]}
}
{
}

\Verse{228}
{
    \zR{贾政叫人去请，那和尚已进来了，也不施礼，也不答话，便往里就跑。}
}
{
    \eR{[English source not present in the checked-in Bencraft/Joly files.]}
}
{
}

\Verse{229}
{
    \zR{贾琏拉 着道：“里头都是内眷，你这野东西混跑什么？”}
}
{
    \eR{[English source not present in the checked-in Bencraft/Joly files.]}
}
{
}

\Verse{230}
{
    \zR{那和尚道：“迟了就不能救了。”}
}
{
    \eR{[English source not present in the checked-in Bencraft/Joly files.]}
}
{
}

\Verse{231}
{
    \zR{贾琏急得一面走，一面乱嚷道：“里头的人不要哭了，和尚进来了！”}
}
{
    \eR{[English source not present in the checked-in Bencraft/Joly files.]}
}
{
}

\Verse{232}
{
    \zR{王夫人等只顾 着哭，那里理会。}
}
{
    \eR{[English source not present in the checked-in Bencraft/Joly files.]}
}
{
}

\Verse{233}
{
    \zR{贾琏走进来又嚷。}
}
{
    \eR{[English source not present in the checked-in Bencraft/Joly files.]}
}
{
}

\Verse{234}
{
    \zR{王夫人等回过头来，见一个长大的和尚，吓了 一跳，躲避不及。}
}
{
    \eR{[English source not present in the checked-in Bencraft/Joly files.]}
}
{
}

\Verse{235}
{
    \zR{那和尚直走到宝玉炕前。}
}
{
    \eR{[English source not present in the checked-in Bencraft/Joly files.]}
}
{
}

\Verse{236}
{
    \zR{宝钗避过一边，袭人见王夫人站着，不 敢走开。}
}
{
    \eR{[English source not present in the checked-in Bencraft/Joly files.]}
}
{
}

\Verse{237}
{
    \zR{只见那和尚道：“施主们，我是送玉来的。”}
}
{
    \eR{[English source not present in the checked-in Bencraft/Joly files.]}
}
{
}

\Verse{238}
{
    \zR{说着，把那块玉擎着道：“快 把银子拿出来，我好救他。”}
}
{
    \eR{[English source not present in the checked-in Bencraft/Joly files.]}
}
{
}

\Verse{239}
{
    \zR{王夫人等惊惶无措，也不择真假，便说道：“若是救活 了人，银子是有的。”}
}
{
    \eR{[English source not present in the checked-in Bencraft/Joly files.]}
}
{
}

\Verse{240}
{
    \zR{那和尚笑道：“拿来！”}
}
{
    \eR{[English source not present in the checked-in Bencraft/Joly files.]}
}
{
}

\Verse{241}
{
    \zR{王夫人道：“你放心，横竖折变的出来。”}
}
{
    \eR{[English source not present in the checked-in Bencraft/Joly files.]}
}
{
}

\Verse{242}
{
    \zR{和尚哈哈大笑，手拿着玉，在宝玉耳边叫道：“宝玉，宝玉!你的‘宝玉’回来了。”}
}
{
    \eR{[English source not present in the checked-in Bencraft/Joly files.]}
}
{
}

\Verse{243}
{
    \zR{说了这一句，王夫人等见宝玉把眼一睁。}
}
{
    \eR{[English source not present in the checked-in Bencraft/Joly files.]}
}
{
}

\Verse{244}
{
    \zR{袭人说道：“好了！”}
}
{
    \eR{[English source not present in the checked-in Bencraft/Joly files.]}
}
{
}

\Verse{245}
{
    \zR{只见宝玉便问道：“在 那里呢？”}
}
{
    \eR{[English source not present in the checked-in Bencraft/Joly files.]}
}
{
}

\Verse{246}
{
    \zR{那和尚把玉递给他手里。}
}
{
    \eR{[English source not present in the checked-in Bencraft/Joly files.]}
}
{
}

\Verse{247}
{
    \zR{宝玉先前紧紧的攥着，后来慢慢的回过手来， 放在自己眼前，细细的一看，说：“嗳呀!久违了。”}
}
{
    \eR{[English source not present in the checked-in Bencraft/Joly files.]}
}
{
}

\Verse{248}
{
    \zR{里外众人都喜欢的念佛，连宝 钗也顾不得有和尚了。}
}
{
    \eR{[English source not present in the checked-in Bencraft/Joly files.]}
}
{
}

\Verse{249}
{
    \zR{贾琏也走过来一看，果见宝玉回过来了，心里一喜，疾忙躲出去了。}
}
{
    \eR{[English source not present in the checked-in Bencraft/Joly files.]}
}
{
}

\Verse{250}
{
    \zR{那和尚也 不言语，赶来拉着贾琏就跑。}
}
{
    \eR{[English source not present in the checked-in Bencraft/Joly files.]}
}
{
}

\Verse{251}
{
    \zR{贾琏只得跟着，到了前头，赶着告诉贾政。}
}
{
    \eR{[English source not present in the checked-in Bencraft/Joly files.]}
}
{
}

\Verse{252}
{
    \zR{贾政听了 喜欢，即找和尚施礼叩谢。}
}
{
    \eR{[English source not present in the checked-in Bencraft/Joly files.]}
}
{
}

\Verse{253}
{
    \zR{和尚还了礼坐下。}
}
{
    \eR{[English source not present in the checked-in Bencraft/Joly files.]}
}
{
}

\Verse{254}
{
    \zR{贾琏心下狐疑：“必是要了银子才走。”}
}
{
    \eR{[English source not present in the checked-in Bencraft/Joly files.]}
}
{
}

\Verse{255}
{
    \zR{贾政细看那和尚，又非前次见的，便问：“宝刹何方？}
}
{
    \eR{[English source not present in the checked-in Bencraft/Joly files.]}
}
{
}

\Verse{256}
{
    \zR{法师大号？}
}
{
    \eR{[English source not present in the checked-in Bencraft/Joly files.]}
}
{
}

\Verse{257}
{
    \zR{这玉是那里得的？}
}
{
    \eR{[English source not present in the checked-in Bencraft/Joly files.]}
}
{
}

\Verse{258}
{
    \zR{怎 么小儿一见便会活过来呢？”}
}
{
    \eR{[English source not present in the checked-in Bencraft/Joly files.]}
}
{
}

\Verse{259}
{
    \zR{那和尚微微笑道：“我也不知道，只要拿一万银子来 就完了。”}
}
{
    \eR{[English source not present in the checked-in Bencraft/Joly files.]}
}
{
}

\Verse{260}
{
    \zR{贾政见这和尚粗鲁，也不敢得罪，便说：“有。”}
}
{
    \eR{[English source not present in the checked-in Bencraft/Joly files.]}
}
{
}

\Verse{261}
{
    \zR{和尚道：“有便快拿来罢， 我要走了。”}
}
{
    \eR{[English source not present in the checked-in Bencraft/Joly files.]}
}
{
}

\Verse{262}
{
    \zR{贾政道：“略请少坐，待我进内瞧瞧。”}
}
{
    \eR{[English source not present in the checked-in Bencraft/Joly files.]}
}
{
}

\Verse{263}
{
    \zR{和尚道：“你去，快出来才好。”}
}
{
    \eR{[English source not present in the checked-in Bencraft/Joly files.]}
}
{
}

\Verse{264}
{
    \zR{贾政果然进去，也不及告诉，便走到宝玉炕前。}
}
{
    \eR{[English source not present in the checked-in Bencraft/Joly files.]}
}
{
}

\Verse{265}
{
    \zR{宝玉见是父亲来，欲要爬起， 因身子虚弱，起不来。}
}
{
    \eR{[English source not present in the checked-in Bencraft/Joly files.]}
}
{
}

\Verse{266}
{
    \zR{王夫人按着说道：“不要动。”}
}
{
    \eR{[English source not present in the checked-in Bencraft/Joly files.]}
}
{
}

\Verse{267}
{
    \zR{宝玉笑着，拿这玉给贾政瞧， 道：“宝玉来了。”}
}
{
    \eR{[English source not present in the checked-in Bencraft/Joly files.]}
}
{
}

\Verse{268}
{
    \zR{贾政略略一看，知道此玉有些根源，也不细看，便和王夫人道： “宝玉好过来了，这赏银怎么样？”}
}
{
    \eR{[English source not present in the checked-in Bencraft/Joly files.]}
}
{
}

\Verse{269}
{
    \zR{王夫人道：“尽着我所有的折变了给他就是了。”}
}
{
    \eR{[English source not present in the checked-in Bencraft/Joly files.]}
}
{
}

\Verse{270}
{
    \zR{宝玉道：“只怕这和尚不是要银子的罢？”}
}
{
    \eR{[English source not present in the checked-in Bencraft/Joly files.]}
}
{
}

\Verse{271}
{
    \zR{贾政点头道：“我也看来古怪，但是他口 口声声的要银子。”}
}
{
    \eR{[English source not present in the checked-in Bencraft/Joly files.]}
}
{
}

\Verse{272}
{
    \zR{王夫人道：“老爷出去先款留着他再说。”}
}
{
    \eR{[English source not present in the checked-in Bencraft/Joly files.]}
}
{
}

\Verse{273}
{
    \zR{贾政出来。}
}
{
    \eR{[English source not present in the checked-in Bencraft/Joly files.]}
}
{
}

\Verse{274}
{
    \zR{宝玉便嚷 饿了，喝了一碗粥，还说要饭。}
}
{
    \eR{[English source not present in the checked-in Bencraft/Joly files.]}
}
{
}

\Verse{275}
{
    \zR{婆子们果然取了饭来。}
}
{
    \eR{[English source not present in the checked-in Bencraft/Joly files.]}
}
{
}

\Verse{276}
{
    \zR{王夫人还不敢给他吃。}
}
{
    \eR{[English source not present in the checked-in Bencraft/Joly files.]}
}
{
}

\Verse{277}
{
    \zR{宝玉 说：“不妨的，我已经好了。”}
}
{
    \eR{[English source not present in the checked-in Bencraft/Joly files.]}
}
{
}

\Verse{278}
{
    \zR{便爬着吃了一碗，渐渐的神气果然好过来了，便要坐 起来。}
}
{
    \eR{[English source not present in the checked-in Bencraft/Joly files.]}
}
{
}

\Verse{279}
{
    \zR{麝月上去轻轻的扶起，因心里喜欢忘了情，说道：“真是宝贝，才看见了一 会儿，就好了。}
}
{
    \eR{[English source not present in the checked-in Bencraft/Joly files.]}
}
{
}

\Verse{280}
{
    \zR{亏的当初没有砸破！”}
}
{
    \eR{[English source not present in the checked-in Bencraft/Joly files.]}
}
{
}

\Verse{281}
{
    \zR{宝玉听了这话，神色一变，把玉一撂，身子 往后一仰。}
}
{
    \eR{[English source not present in the checked-in Bencraft/Joly files.]}
}
{
}

\Verse{282}
{
    \zR{未知死活，下回分解。}
}
{
    \eR{[English source not present in the checked-in Bencraft/Joly files.]}
}
{
}

\Verse{283}
{
    \zR{(本章完)}
}
{
    \eR{[English source not present in the checked-in Bencraft/Joly files.]}
}
{
}
